\begin{ex}
Số 24 568 là số chẵn hay số lẻ? (Nếu là số chẵn ghi 0, số lẻ ghi 1).
\shortans{0}
\loigiai{
Chữ số tận cùng là 8 nên là số chẵn (0).
}
\end{ex}

\begin{ex}
Số 73 491 là số chẵn hay số lẻ? (Nếu là số chẵn ghi 0, số lẻ ghi 1).
\shortans{1}
\loigiai{
Chữ số tận cùng là 1 nên là số lẻ (1).
}
\end{ex}

\begin{ex}
Trong các số $124, 357, 460, 589$, có bao nhiêu số chẵn?
\shortans{2}
\loigiai{
Các số chẵn là 124, 460 (có 2 số).
}
\end{ex}

\begin{ex}
Trong các số $135, 246, 579, 681, 702$, có bao nhiêu số lẻ?
\shortans{3}
\loigiai{
Các số lẻ là 135, 579, 681 (có 3 số).
}
\end{ex}

\begin{ex}
Số chẵn liền sau của số 48 là số nào?
\shortans{50}
\loigiai{
$48 + 2 = 50$.
}
\end{ex}

\begin{ex}
Số chẵn liền trước của số 100 là số nào?
\shortans{98}
\loigiai{
$100 - 2 = 98$.
}
\end{ex}

\begin{ex}
Số lẻ liền sau của số 79 là số nào?
\shortans{81}
\loigiai{
$79 + 2 = 81$.
}
\end{ex}

\begin{ex}
Số lẻ liền trước của số 51 là số nào?
\shortans{49}
\loigiai{
$51 - 2 = 49$.
}
\end{ex}

\begin{ex}
Số chẵn bé nhất có hai chữ số là số nào?
\shortans{10}
\loigiai{
Số đó là 10.
}
\end{ex}

\begin{ex}
Số lẻ bé nhất có hai chữ số là số nào?
\shortans{11}
\loigiai{
Số đó là 11.
}
\end{ex}

\begin{ex}
Số chẵn lớn nhất có hai chữ số là số nào?
\shortans{98}
\loigiai{
Số đó là 98.
}
\end{ex}

\begin{ex}
Số lẻ lớn nhất có hai chữ số là số nào?
\shortans{99}
\loigiai{
Số đó là 99.
}
\end{ex}

\begin{ex}
Số chẵn bé nhất có ba chữ số là số nào?
\shortans{100}
\loigiai{
Số đó là 100.
}
\end{ex}

\begin{ex}
Số lẻ bé nhất có ba chữ số là số nào?
\shortans{101}
\loigiai{
Số đó là 101.
}
\end{ex}

\begin{ex}
Số chẵn lớn nhất có ba chữ số khác nhau là số nào?
\shortans{986}
\loigiai{
Số đó là 986.
}
\end{ex}

\begin{ex}
Số lẻ bé nhất có ba chữ số khác nhau là số nào?
\shortans{103}
\loigiai{
Số đó là 103.
}
\end{ex}

\begin{ex}
Tổng của hai số chẵn là số chẵn hay số lẻ? (Chẵn ghi 0, lẻ ghi 1).
\shortans{0}
\loigiai{
Tổng hai số chẵn luôn là số chẵn (0).
}
\end{ex}

\begin{ex}
Tổng của hai số lẻ là số chẵn hay số lẻ? (Chẵn ghi 0, lẻ ghi 1).
\shortans{0}
\loigiai{
Tổng hai số lẻ luôn là số chẵn (0).
}
\end{ex}

\begin{ex}
Tổng của một số chẵn và một số lẻ là số chẵn hay lẻ? (Chẵn ghi 0, lẻ ghi 1).
\shortans{1}
\loigiai{
Tổng của một số chẵn và một số lẻ luôn là số lẻ (1).
}
\end{ex}

\begin{ex}
Tích của hai số lẻ là số chẵn hay lẻ? (Chẵn ghi 0, lẻ ghi 1).
\shortans{1}
\loigiai{
Tích hai số lẻ luôn là số lẻ (1).
}
\end{ex}

\begin{ex}
Tích của một số chẵn với bất kỳ số tự nhiên nào là số chẵn hay lẻ? (Chẵn ghi 0, lẻ ghi 1).
\shortans{0}
\loigiai{
Tích luôn là số chẵn (0).
}
\end{ex}

\begin{ex}
Từ 1 đến 20 có bao nhiêu số chẵn?
\shortans{10}
\loigiai{
Các số chẵn: 2, 4, ..., 20 (có 10 số).
}
\end{ex}

\begin{ex}
Từ 1 đến 20 có bao nhiêu số lẻ?
\shortans{10}
\loigiai{
Các số lẻ: 1, 3, ..., 19 (có 10 số).
}
\end{ex}

\begin{ex}
Hiệu của hai số lẻ là số chẵn hay lẻ? (Chẵn ghi 0, lẻ ghi 1).
\shortans{0}
\loigiai{
Hiệu hai số lẻ là số chẵn (0).
}
\end{ex}

\begin{ex}
Có bao nhiêu số chẵn có hai chữ số?
\shortans{45}
\loigiai{
Có $(98 - 10) : 2 + 1 = 45$ số.
}
\end{ex}
